\section{Sub-Gaussian labels and an expected-risk corollary}
\label{sec:general-shallow-subgaussian-expectation}
The finite-variance confidence theorem does not imply an
expected-risk theorem. For marginally sub-Gaussian labels, however,
the sharper logarithmic-confidence bounded-label argument gives both.
Assume $\mathbb E\exp(Y^2/K_Y^2)\le2$, with a fixed known envelope.
Let $\lambda'_n$ be (GS1) with confidence $\delta/4$, and put
\[
 M_n=\max\{1,K_Y\sqrt{c\log((n+2)/\delta)}\},\qquad
 \lambda_n=M_n^2\lambda'_n,                              \tag{GG1}
\]
for a sufficiently large fixed numerical $c$.
The actual likelihood still uses the original labels and unbounded
raw output, with variance $n\lambda_n\nu/N$.

\begin{corollary}[Marginal sub-Gaussian confidence and expectation]
At every $N\ge n$, $0<\nu\le1$, and every sample size with
$\lambda'_n\le1$, the sharp path and coefficient-budget oracles
(GS5)--(GS6) hold with the same polynomial powers and additional
logarithmic factors, for the raw posterior-average and mean risks.
Their confidence is at least $1-\delta$.
Choosing $\delta=n^{-20}$ in (GG1) also gives an expected-risk
oracle with the same powers. In particular, at fixed dimension and
coefficient cap, every declared width-$C$ comparison class has
\[
 \mathbb E_S R(\bar z_{S})\le
 \mathbb E_S\Pi_S R(z)
 \le \inf_{g\in\mathcal G_{C,B}}R(g)
          +\widetilde O((C+1)B^2/\sqrt n).
\]
The prior, prediction convention, and all observations are unchanged.
This expectation statement is specific to the present sub-Gaussian
label hypothesis; it is not asserted under only a second moment.
\end{corollary}
\begin{proof}
Use the proof-only label $Y_M=\operatorname{clip}(Y,[-M_n,M_n])$.
Except with probability at most $\delta/4$, every observed label
already equals its clipped version. Analyze the observable $z/M_n$
and label $Y_M/M_n$, keeping exactly the original parameter prior.
The likelihood then has normalized scale $\lambda'_n$.
Its prior empirical risk is at most two, and its parameter energy
is still the same chi-square quadratic form.

For a physical comparator $g$, the minimum-mass LP compresses the
physical midpoint $(z+g)/2$ as before, with center energy at most
$P(z)+P_g$. Dividing loss by $M_n^2$ leaves (GM1) valid with
scale $\lambda'_n$ and the same physical path penalty.
The normalized function has path mass $P(z)/M_n\le P(z)$;
the normalized comparator has mass $P_g/M_n\le P_g$.
Consequently the atom, relative-square, and multiplier estimates
used in (GS5) are all valid with these larger physical path bounds.
This proves the normalized oracle. Multiplication by $M_n^2$
preserves its comparison coefficient one and adds only logarithms.

This also transfers population risks. On the clipped-label sample,
\[
 \Pi Q^2\le C[d+1+(\lambda'_n)^{-1}]^2,\qquad
 |z(x)|\le Q/\sqrt2.
\]
Therefore the difference between the posterior's risks under $Y$
and $Y_M$ is at most
\[
 C[d+1+(\lambda'_n)^{-1}]
       \|Y-Y_M\|_{L^2}+\mathbb E|Y^2-Y_M^2|.
\]
It is negligible for sufficiently large $c$ in (GG1).
Comparisons with $R(g)>R(0)$ use the zero comparator at the same
schedule; otherwise $\|g\|_{L^2}\le2\|Y\|_{L^2}$.
For a fixed comparison budget choose a deterministic population minimizer with
$R(g)\le R(0)$; then $\|g\|_{L^2}\le2\|Y\|_{L^2}$.
Its analogous population-risk transfer is negligible by the same
estimate. No input distribution is conditioned on a label event.

It remains to control the exceptional samples for the expectation
claim. For every sample, without any label bound, the Gaussian
evidence and energy estimate give
\[
 \Pi Q^2\le
 C\left[d+1+\frac{1+\mathbb P_nY^2}{\lambda_n}\right]^2.
\]
Thus, writing $M_2=\mathbb EY^2$,
\[
 \Pi R(z)\le C\left[M_2+(d+1)^2+
                \frac{(1+\mathbb P_nY^2)^2}{\lambda_n^2}\right].
\]
The sub-Gaussian assumption gives
$\mathbb E(\mathbb P_nY^2)^4\le\mathbb EY^8<\infty$.
Hence the $L^2$ norm, over the data, of this risk upper bound is at
most $C_{K_Y}[1+d^2+\lambda_n^{-2}]$.
At $\delta=n^{-20}$, Cauchy--Schwarz bounds the expected risk on
the exceptional event by
$C_{K_Y}[1+d^2+\lambda_n^{-2}]n^{-10}$.
In the stated sample-size range, $d\le Cn$ and
$\lambda_n\ge n^{-2}$, so this contribution is $O_{K_Y}(n^{-6})$.
Integrating the confidence bound proves the expected oracle.
\end{proof}

\paragraph{The bounded-monotone companion.}
The same confidence-to-expectation step applies to the paper's
bounded-monotone-activation midpoint oracle of
Theorem~\ref{thm:uncapped-all-d}, with its condition $N\lambda\ge c_\phi$.
This paragraph does not broaden the all-width finite-variance theorem
(AW3) to expectation. Its deterministic
readout-energy bound has the smaller form
$\Pi R(z)\le C_\phi[M_2+1+(1+\mathbb P_nY^2)/\lambda]$.
Thus its already stated logarithmic-confidence sub-Gaussian oracle
also yields an expected-risk version by taking $\delta=n^{-20}$;
the fourth label moment suffices for this final integration step.
The finite-second-moment confidence-only conclusions remain separate.
